\subsection{Irreducible representations}

Lemma 1. --- Every irreducible unitary representation of G is square-integrable.

Indeed, the matrix coefficients of an irreducible unitary representation are continuous and bounded, and are therefore square-integrable on G, since G is compact.

PROPOSITION 2. --- Let π be an irreducible unitary representation of G in a Hilbert space E. The dimension of E is finite and equal to the formal degree of π.

The representation \(\pi\) being square-integrable (Lemma 1), its formal degree \(c_{\mu}(\pi)\) relative to \(\mu\) is defined (def. 4 of V, p. 423); it is a strictly positive real number. Let \((e_i)_{i \in I}\) be a finite orthonormal family in E. Let \(x\) be an element of E with norm 1 (such exists since E is nonzero). For \(i \in I\) , we have

\[
c _ {\mu} (\pi) \int_ {\mathrm{G}} | \langle e _ {i} | \pi (g) x \rangle | ^ {2} d \mu (g) = 1
\]

(prop. 8 of V, p. 424). Summing this formula over \(i \in I\) , we obtain

\[
\begin{array}{r l} & {\mathrm{Card(I)} = c _ {\mu} (\pi) \int_ {\mathrm{G}} \sum_ {i \in \mathrm{I}} | \langle e _ {i} | \pi (g) x \rangle | ^ {2} d \mu (g)} \\ & {\qquad \leqslant c _ {\mu} (\pi) \int_ {\mathrm{G}} \| \pi (g) x \| ^ {2} d \mu (g) = c _ {\mu} (\pi)} \end{array}
\]

by Bessel's inequality (EVT, V, p. 21, prop. 4). Hence the cardinality of I is bounded above by \(c_{\mu}(\pi)\) . This implies that the dimension of E is finite.

One can then apply the preceding result to an orthonormal basis \((e_{i})_{i\in\mathrm{I}}\) of E. We obtain, by EVT, V, p. 22, prop. 5,

\[
\begin{array}{r} \dim (\mathrm{E}) = c _ {\mu} (\pi) \int_ {\mathrm{G}} \sum_ {i \in \mathrm{I}} | \langle e _ {i} | \pi (g) x \rangle | ^ {2} d \mu (g) \\ = c _ {\mu} (\pi) \int_ {\mathrm{G}} \| \pi (g) x \| ^ {2} d \mu (g) = c _ {\mu} (\pi), \end{array}
\]

which concludes the proof.

In particular, it is therefore permissible to speak of the character \(\chi_{\pi}\) of an irreducible unitary representation \(\pi\) of G. It is a continuous function on G.

COROLLARY 1. --- Let \(\pi\) be an irreducible unitary representation of G in a Hilbert space E. The character of \(\pi\) is a Hermitian element of the involutive algebra \(\mathrm{L}^{1}(\mathrm{G})\) .

Let \(x \in G\) . Since \(G\) is unimodular, we have \(\chi_{\pi}^{*}(x) = \overline{\operatorname{Tr}(\pi(x^{-1}))}\) , whence \(\chi_{\pi}^{*}(x) = \operatorname{Tr}(\pi(x))\) since \(\pi\) is a unitary representation.

COROLLARY 2. --- a) Let \(\pi\) be an irreducible unitary representation of G in a Hilbert space E. We have

\[
\int_ {\mathrm{G}} \overline {{\langle x \mid \pi (g) x ^ {\prime} \rangle}} \langle y \mid \pi (g) y ^ {\prime} \rangle d \mu (g) = \frac {1}{\dim (\mathrm{E})} \overline {{\langle x \mid y \rangle}} \langle x ^ {\prime} \mid y ^ {\prime} \rangle
\]

for all \((x,y,x',y')\in \mathrm{E}^4\) .

\begin{enumerate}
\def\labelenumi{\alph{enumi})}
\setcounter{enumi}{1}
\tightlist
\item
  Let \(\pi_{1}\) (resp. \(\pi_{2}\) ) be an irreducible unitary representation of G in a Hilbert space \(E_{1}\) (resp. \(E_{2}\) ). If \(\pi_{1}\) and \(\pi_{2}\) are not isomorphic, one has
\end{enumerate}

\[
\int_ {G} \overline {{\langle x \mid \pi_ {1} (g) x ^ {\prime} \rangle}} \langle y \mid \pi_ {2} (g) y ^ {\prime} \rangle d \mu (g) = 0
\]

for all \((x,x^{\prime},y,y^{\prime})\in \mathrm{E}_{1}^{2}\times \mathrm{E}_{2}^{2}\)

This follows immediately from prop. 8 of V, p. 424 and prop. 9 of V, p. 424, taking into account lemma 1 and the formula \(c_{\mu}(\pi) = \dim(\mathrm{E})\) (prop. 2).

COROLLARY 3. --- Let \(\pi_1\) and \(\pi_2\) be irreducible representations of G. In the Hilbert space \(\mathrm{L}^{2}(\mathrm{G})\) , one has \(\langle \chi_{\pi_1} | \chi_{\pi_2} \rangle = 1\) if \(\pi_1\) and \(\pi_2\) are isomorphic and \(\langle \chi_{\pi_1} | \chi_{\pi_2} \rangle = 0\) otherwise. In other words, the family of characters of the classes \(\pi \in \widehat{\mathbf{G}}\) is an orthonormal family in \(\mathrm{L}^{2}(\mathrm{G})\) .

This follows from proposition 1 of V, p. 457 and corollary 2, noting that for every orthonormal basis \((e_{i})_{i\in\mathrm{I}}\) of the space of a finite-dimensional unitary representation \(\pi\) of G, and for every \(g\in G\) , we have the formula

\[
\chi_ {\pi} (g) = \sum_ {i \in I} \langle e _ {i} | \pi (g) e _ {i} \rangle .
\]

COROLLARY 4. --- a) Let \(\varrho\) be a continuous finite-dimensional representation of G. One has

\[
\chi_ {\varrho} = \sum_ {\pi \in \widehat {G}} \dim (\operatorname{Hom} _ {G} (\pi , \varrho)) \chi_ {\pi} = \sum_ {\pi \in \widehat {G}} \dim (\operatorname{Hom} _ {G} (\varrho , \pi)) \chi_ {\pi}.
\]

\begin{enumerate}
\def\labelenumi{\alph{enumi})}
\setcounter{enumi}{1}
\tightlist
\item
  Let \(\varrho_{1}\) and \(\varrho_{2}\) be continuous finite-dimensional representations of G. One has
\end{enumerate}

\[
\left\langle \chi_ {\varrho_ {1}} \mid \chi_ {\varrho_ {2}} \right\rangle = \dim (\operatorname{Hom} _ {\mathrm{G}} (\varrho_ {1}, \varrho_ {2})) = \dim (\operatorname{Hom} _ {\mathrm{G}} (\varrho_ {2}, \varrho_ {1})).
\]

\begin{enumerate}
\def\labelenumi{\alph{enumi})}
\setcounter{enumi}{2}
\tightlist
\item
  A continuous finite-dimensional representation \(\varrho\) of G is irreducible if and only if \(\| \chi_{\varrho}\|^{2} = 1\) .
\end{enumerate}

Since \(\varrho\) is semi-simple, it is isomorphic to the direct sum of its \(\pi\)-isotypic components \(\mathrm{M}_{\pi}(\varrho)\) for \(\pi \in \widehat{\mathrm{G}}\) . Let \(m_{\pi}(\varrho)\) denote the multiplicity of \(\pi\) in \(\varrho\) (def. 10 of V, p. 398); we then have

\[
\chi_ {\varrho} = \sum_ {\pi \in \widehat {G}} m _ {\pi} (\varrho) \chi_ {\pi}.
\]

The assertion a) therefore follows from the fact that

\[
m _ {\pi} (\varrho) = \dim \operatorname{Hom} _ {\mathrm{G}} (\pi , \varrho) = \dim \operatorname{Hom} _ {\mathrm{G}} (\varrho , \pi)
\]

(formula (1) of V, p. 377 and cor. 2 of V, p. 387).

By bilinearity and orthonormality of characters, the assertion \(a\) ) implies that

\[
\langle \chi_ {\varrho_ {1}} \mid \chi_ {\varrho_ {2}} \rangle = \sum_ {\pi \in \widehat {\mathrm{G}}} \dim (\operatorname{Hom} _ {\mathrm{G}} (\pi , \varrho_ {1})) \dim (\operatorname{Hom} _ {\mathrm{G}} (\pi , \varrho_ {2})),
\]

and on the other hand we have canonical isomorphisms

\[
\begin{array}{c} \operatorname{Hom} _ {\mathrm{G}} (\varrho _ {1}, \varrho _ {2}) \to \bigoplus_ {(\pi _ {1}, \pi _ {2}) \in \widehat {\mathrm{G}} \times \widehat {\mathrm{G}}} \operatorname{Hom} _ {\mathrm{G}} (\mathrm{M} _ {\pi _ {1}} (\varrho _ {1}), \mathrm{M} _ {\pi _ {2}} (\varrho _ {2})) \\ \qquad \qquad \qquad \qquad \qquad \qquad \qquad \qquad \qquad \qquad \qquad \qquad \qquad \qquad \qquad \qquad \qquad \qquad \qquad \qquad \qquad \qquad \qquad \qquad \qquad \qquad \qquad \qquad \qquad \qquad \qquad \qquad \qquad \qquad \qquad \to \bigoplus_ {\pi \in \widehat {\mathrm{G}}} \operatorname{Hom} _ {\mathrm{G}} (\mathrm{M} _ {\pi} (\varrho _ {1}), \mathrm{M} _ {\pi} (\varrho _ {2})) \end{array}
\]

(formula (1) of V, p. 377), whence it follows that

\[
\dim (\operatorname{Hom} _ {\mathrm{G}} (\varrho_ {1}, \varrho_ {2})) = \sum_ {\pi \in \widehat {\mathrm{G}}} m _ {\pi} (\varrho_ {1}) m _ {\pi} (\varrho_ {2}),
\]

whence the assertion \(b\) ). The assertion \(c\) ) also follows from \(a\) ) and Schur's lemma (prop. 6 of V, p. 386).

COROLLARY 5. --- Let \(\pi_1\) and \(\pi_2\) be irreducible unitary representations of G. Then \(\pi_1(\overline{\chi}_{\pi_2}) = 0\) if \(\pi_1\) is not isomorphic to \(\pi_2\) , and \(\pi_1(\overline{\chi}_{\pi_1})\) is multiplication by \(1/\dim(\pi_1)\) .

The character of \(\pi_2\) is a continuous central function on G, hence the linear map \(u = \pi_1(\overline{\chi}_{\pi_2})\) is defined and belongs to the space \(\mathrm{Hom}_{\mathrm{G}}(\pi_1,\pi_1)\) . It is a homothety by Schur's lemma (prop. 6 of V, p. 386) whose trace is

\[
\begin{array}{c} \operatorname{Tr} (u) = \operatorname{Tr} \Bigl (\int_ {\mathrm{G}} \overline {{\chi_ {\pi_ {2}} (g)}}   \pi_ {1} (g) d \mu (g) \Bigr) \\ = \int_ {\mathrm{G}} \overline {{\chi_ {\pi_ {2}} (g)}} \chi_ {\pi_ {1}} (g) d \mu (g). \end{array}
\]

By the orthogonality relations, the trace of \(u\) is therefore zero if \(\pi_1\) is not isomorphic to \(\pi_2\) and equals 1 otherwise. The assertion follows.

Remark. --- When G is finite, the Schur orthogonality relation (resp. the character orthogonality formula) coincides with that of A, VIII, p. 399 (resp. that of A, VIII, p. 400, prop. 4). On the other hand, the "second orthogonality relation" of characters of finite groups (A, VIII, p. 402, formula (32)) has no exact analogue when G is compact.

For G finite, the particular case of the second orthogonality formula corresponding to the conjugacy class of e is the formula

\[
\frac {1}{\mathrm{CardG}} \sum_ {\pi \in \widehat {\mathrm{G}}} \dim (\pi) \chi _ {\pi} (g) = \left\{ \begin{array}{l l} 1 & \text {if} g = e \\ 0 & \text {otherwise}, \end{array} \right.
\]

which is also interpreted as the computation of the character of the regular representation \(\gamma_{G}\) of G, and which is equivalent to the formula \(\mathrm{Tr}(\boldsymbol{\gamma}_{\mathrm{G}}(f)) = f(e)\) for every function f from G to C.

Let G again be an arbitrary compact group. For every function \(f \in \mathcal{C}(\mathrm{G})\) , the endomorphism \(\boldsymbol{\gamma}_{\mathrm{G}}(f)\) of \(\mathrm{L}^{2}(\mathrm{G})\) coincides with the endomorphism \(\varphi \mapsto f * \varphi\) and it has finite trace (lemma 4 of V, p. 407 and corollary 2 of III, p. 33). By loc. cit. and th. 2 of IV, p. 177, we also have

\[
\operatorname{Tr} \left(\boldsymbol {\gamma} _ {\mathrm{G}} (f)\right) = \int_ {\mathrm{G}} f \left(x ^ {- 1} x\right) d \mu (x) = f (e).
\]

PROPOSITION 3. --- Suppose that \(G = G_{1} \times G_{2}\) where \(G_{1}\) and \(G_{2}\) are compact groups. The map b from \(\widehat{G}_{1} \times \widehat{G}_{2}\) to \(\widehat{G}\) that sends \((\pi_{1}, \pi_{2})\) to the class of \(\pi_{1} \boxtimes \pi_{2}\) is a bijection.

By cor. 6 of V, p. 389, the map b is well defined.

Let \(\pi_1\) and \(\pi_2\) be elements of \(\widehat{\mathbf{G}}\) . Let \(\pi \in \widehat{\mathbf{G}}\) . The \(\pi\)-isotypic component of the restriction \(\varpi\) of \(\pi_1 \boxtimes \pi_2\) to \(\mathrm{G}_1 \times \{e\}\) is equal to \(\varpi\) if \(\pi_1 = \pi\) and is zero otherwise (lemma 8 of V, p. 384). Thus, the unitary representation \(\pi_1 \boxtimes \pi_2\) determines \(\pi_1\) up to isomorphism; likewise, it determines \(\pi_2\) , which proves that \(b\) is injective.

Let us prove that \(b\) is surjective. Let \(\pi\) be an irreducible unitary representation of \(G_1 \times G_2\) in a Hilbert space E. It is finite-dimensional (prop. 2). Let \(\pi_1\) be an irreducible unitary representation of \(G_1\) in a Hilbert space \(E_1\) such that the restriction of \(\pi\) to \(G_1 \times \{e\}\) contains a subrepresentation isomorphic to \(\pi_1\) (lemma 3 of V, p. 379). Let \(F = \text{Hom}_{G_1}(\pi_1, \pi)\) . This is a nonzero finite-dimensional vector space. For \(h \in G_2\) and \(u \in F\) , let \(\varrho(h)(u)\) be the linear map from \(E_1\) to E defined by \(x \mapsto \pi(e, h)(u(x))\) . Since \(G_1 \times \{e\}\) and \(\{e\} \times G_2\) commute in G, the map \(\varrho(h)(u)\) belongs to the space F. The map \(\varrho\) is a linear representation of \(G_2\) in the space F; it is continuous. Since F is nonzero, there exists an irreducible subrepresentation \(\pi_2\) of \(\varrho\) (lemma 3 of V, p. 379). Let \(E_2\) be the space of \(\pi_2\) . The linear map \(v: E_1 \otimes E_2 \to E\) such that \(x \otimes u \mapsto u(x)\) for \(x \in E_1\) and \(u \in E_2\) is then a nonzero G-morphism from \(\pi_1 \boxtimes \pi_2\) to \(\pi\); since the representations \(\pi\) and \(\pi_1 \boxtimes \pi_2\) are irreducible and finite-dimensional, the morphism \(v\) is an isomorphism (Lemma 2 of V, p. 378).

This proposition is to be compared with Theorem 1 of A, VIII, p. 208.

